% généré par scripts/06_textes.py
\dosmeta{Année de classement : 369 (directe : nom consulaire lu). Références : ICUR II 4165 = ILCV 1603 (add, em). Rattachement selon l’EDB : \emph{Basilica S.Petri apostoli}. Support : Support non précisé par l’EDB. Reproductions signalées par l’EDB : dessin, ICVR I (de Rossi) (droits réservés ; non reproduites). Fiche EDB : Antonio Enrico Felle, première version 2003. Notice du catalogue p.~\pageref{not:133}.}
\txtbloc{Texte transmis (EDCS, état 2022)}{\lnum{1}Felix sanctae fidei vocitus iit in pace cuius tantus amor et caritas retenetur ab amicis in aevo\par
\lnum{2}qui cum esset fuit solacius misericors omnibus notus\par
\lnum{3}Agrippina fecit dulcissimo suo marito\par
\lnum{4}cum qu⸢o⸣ vixit sine l(a)esione animi annos III et m(enses) X\par
\lnum{5}fuit in saeculum quo(a)d vixit annos XXXII dep(ositus) XIII Kal(endas) Sept(embres) Valentiniano\par
\lnum{6}n(obilissimo) p(uero) et Victori conss(ulibus)}
\txtbloc{Texte de l’EDB (EDB19224)}{\lnum{1}Felix sanctae fidei vocitus iit in pace\par
\lnum{2}cuius tantus amor et caritas retenetur ab amicis in aevo\par
\lnum{3}qui cum esset fuit solacius misericors omnibus notus\par
\lnum{4}Agrippina fecit dulcissimo suo marito\par
\lnum{5}cum quem vixit sine lesione animi annos III det m(enses) X\par
\lnum{6}fuit in saeculum quod vixit annos XXXII dep(ositus) XIII kal(endas) sept(embres) Valentiniano n(obilissimo) p(uero) et Victori conss(ulibus)}
